% =====================================================================
% 07-challenges.tex
% =====================================================================
\section{Open Challenges and Future Directions}\label{sec:challenges}

The field's open problems fall into three broad axes: reliable reward design, efficient optimization and long-horizon scalability. These axes give rise to the specific challenges  discussed in the subsections below.

\subsection{Reward Hacking and Robust Alignment}

Learned the reward proxies remain vulnerable to exploitation, often favoring verbosity, formatting or sycophancy over genuine quality. 
Overoptimization follows predictable scaling laws~\cite{gao2023scaling}, and reward-model ensembles mitigate but do not eliminate these failures~\cite{coste2024ensembles}. Alignment is also fragile, cheap fine-tuning can undo safety training~\cite{qi2024finetuning,lermen2024lora}.
Verifiable rewards avoid many of these issues when a suiteble verifier exists. Current directions include uncertainty-calibrated reward
models, online reward-hacking detection, and information-theoretic

\subsection{Credit Assignment and Sparse Rewards}

In multi-step settings, the central difficulty is attributing a delayed,
sparse outcome to individual actions in a long trajectory. Process reward
models~\cite{lightman2023verify,wang2024omegaprm} provide denser signal,
but generating meaningful intermediate rewards without extensive
annotation, and beyond structured mathematics, remains open.
Trajectory-aware advantage estimation and automated sub-goal verification
are active research avenues.

\subsection{Sample Efficiency and Cost}

Policy optimization remains expensive under online RL becouse each update requires sampling from the full model.
Asynchronous, off-policy RLHF~\cite{noukhovitch2025async}, difficulty-targeted data selection, and rollout reuse with importance sampling all aim to reduce the number of rollouts. Bridging offline pre-training of policies with a short online
phase offers an attractive compromise between cost and performance.

\subsection{Scalable Oversight}

As models approach or exceed human ability on a task, human preference
ceases to be a reliable reward signal, since annotators cannot judge what they
cannot verify~\cite{casper2023open}. Scalable-oversight proposals include
debate, recursive task decomposition, and weak-to-strong generalization.
The broad migration toward verifiable rewards is, in part, a response to
this limitation, but most real tasks lack a crisp verifier.

\subsection{Evaluation and Benchmarks}

Evaluation has not kept pace with model capabilities. Most benchmarks focus on single-turn interactions,
while the behaviors of interest arise in is multi-turn, adaptive settings. Benchmarks are also prone to contamination as models scale. The field needs
standardized, contamination-resistant suites that measure long-horizon, tool-using competence and the robustness of alignment under distribution
shift.
